% =============================================================================
%  §4  Ch.7 位移插值（Displacement Interpolation）
%  来源：Villani, Optimal Transport: Old and New, Chapter 7 (pp. 114–160)
%  主要定理：定理 7.21
% =============================================================================
\TLsection{4 \ 位移插值}

% ── 动机（一）：线性插值的缺陷 ──────────────────────────────────────────────
\begin{frame}{动机：线性插值的缺陷}
  \textbf{自然问题。} 给定两个概率分布 $\mu_{0}$ 和 $\mu_{1}$，如何构造"合理的"中间分布？

  \textbf{线性插值：} $\mu_{t}^{\mathrm{lin}}=(1-t)\mu_{0}+t\mu_{1}$
  \begin{itemize}
    \item 数学上简单，但\emph{物理上不合理}：质量被"传送"（teleport）而非移动。
    \item $\mu_{0}=\delta_{0}$，$\mu_{1}=\delta_{1}$ 时，$\mu_{1/2}^{\mathrm{lin}}=\tfrac{1}{2}\delta_{0}+\tfrac{1}{2}\delta_{1}$——质量\emph{分裂}到两处。
    \item 两高斯分布的线性插值在 $t=\tfrac{1}{2}$ 产生\emph{双峰}（bimodal），与直觉相悖。
  \end{itemize}

  \textbf{位移插值（McCann 1995）：} 让每单位质量沿\TLterm{最优路径}从 $\mu_{0}$ 移动到 $\mu_{1}$；
  $\mu_{t}$ 是所有质量在 $t$ 时刻位置的分布。

  \TLtakeaway{线性插值"传送"（teleport）质量产生双峰；位移插值沿最优路径整体平移，保留分布形状。}
\end{frame}

% ── 动机（图示）：线性 vs 位移插值 ───────────────────────────────────────────
\begin{frame}{图示：线性插值 vs 位移插值（对比）}
  \begin{columns}[T]
    \begin{column}{0.48\textwidth}
      \centering
      \textbf{线性插值}（产生双峰）\\[4pt]
      \begin{tikzpicture}[scale=0.72,>={Stealth[length=3pt]}]
        % t=0：单峰在 x=0（与位移列同一端点）
        \draw[TLprimary,thick,domain=-1.0:1.0,smooth,samples=25] plot(\x,{1.0*exp(-3*\x*\x)});
        \node[TLprimary,font=\scriptsize,below] at (0,-0.1) {$t=0$};
        \begin{scope}[shift={(0,-1.8)}]
          % t=1/2：½在 0 + ½在 1.6，即两端点的平均 → 双峰
          \draw[TLprimary!50!TLaccent,thick,domain=-1.0:2.6,smooth,samples=50]
            plot(\x,{0.5*exp(-3*\x*\x)+0.5*exp(-3*(\x-1.6)*(\x-1.6))});
          \node[TLprimary!50!TLaccent,font=\scriptsize,below] at (0.8,-0.1) {$t=\tfrac12$（双峰）};
        \end{scope}
        \begin{scope}[shift={(0,-3.6)}]
          % t=1：单峰在 x=1.6（与位移列同一端点）
          \draw[TLaccent,thick,domain=0.6:2.6,smooth,samples=25] plot(\x,{1.0*exp(-3*(\x-1.6)*(\x-1.6))});
          \node[TLaccent,font=\scriptsize,below] at (1.6,-0.1) {$t=1$};
        \end{scope}
      \end{tikzpicture}
    \end{column}
    \begin{column}{0.48\textwidth}
      \centering
      \textbf{位移插值}（平滑平移）\\[4pt]
      \begin{tikzpicture}[scale=0.72,>={Stealth[length=3pt]}]
        \draw[TLprimary,thick,domain=-1.0:1.0,smooth,samples=25]
          plot(\x,{1.0*exp(-3*\x*\x)});
        \node[TLprimary,font=\scriptsize,below] at (0,-0.1) {$t=0$};
        \begin{scope}[shift={(0,-1.8)}]
          \draw[TLprimary!50!TLaccent,thick,domain=-0.2:1.8,smooth,samples=25]
            plot(\x,{1.0*exp(-3*(\x-0.8)*(\x-0.8))});
          \node[TLprimary!50!TLaccent,font=\scriptsize,below] at (0.8,-0.1) {$t=\tfrac12$（单峰）};
        \end{scope}
        \begin{scope}[shift={(0,-3.6)}]
          \draw[TLaccent,thick,domain=0.6:2.6,smooth,samples=25]
            plot(\x,{1.0*exp(-3*(\x-1.6)*(\x-1.6))});
          \node[TLaccent,font=\scriptsize,below] at (1.6,-0.1) {$t=1$};
        \end{scope}
      \end{tikzpicture}
    \end{column}
  \end{columns}
  \vspace{2pt}
  {\scriptsize 据 Villani 重绘。}

  \TLtakeaway{位移插值中单峰始终保持单峰并平滑平移；线性插值在中间产生双峰——这正是最优传输几何直觉的来源。}
\end{frame}

% ── 作用量泛函与 Lagrange 量 ─────────────────────────────────────────────────
\begin{frame}{作用量泛函与 Lagrange 量（7.2）}
  \deflab{作用量泛函（Villani 第 7 章）。}
  \deflab{补一个事实：切丛。} $TM$ 把位置 $x\in M$ 与速度 $v\in T_{x}M$ 配成点 $(x,v)$。
  \deflab{补一个事实：Lagrange 量。} $L(x,v,t)$ 表示时刻 $t$ 以速度 $v$ 经过 $x$ 的瞬时运动成本。
  \TLterm{作用量}把瞬时成本沿曲线累加：
  \[
    A(\gamma) = \int_{0}^{1} L(\gamma_{t},\dot{\gamma}_{t},t)\,\dd t.
  \]
  \textbf{典型形式：} $L=\tfrac{|v|^{2}}{2}$，$A(\gamma)=\int_{0}^{1}\tfrac{|\dot{\gamma}_{t}|^{2}}{2}\dd t$（动能积分）。

  \textbf{最小化曲线与成本：} 同端点下 $A(\gamma)\le A(\tilde{\gamma})$；关联成本
  $c(x,y):=\inf_{\gamma_{0}=x,\,\gamma_{1}=y}A(\gamma)$。
  \textbf{例 7.1。} $L=|v|$ 时 $A(\gamma)=\mathrm{length}(\gamma)$，$c(x,y)=|x-y|$。

  \TLtakeaway{作用量 = Lagrange 量的时间积分；最小化对应"最省力"路径；关联成本 $c(x,y)$ 是最优作用量值。}
\end{frame}

% ── 例 7.2 与例 7.4 ──────────────────────────────────────────────────────────
\begin{frame}{例 7.2 \& 7.4：凸 Lagrange 量与 Riemann 测地线}
  \deflab{例 7.2（Villani Example 7.2）。} 取 $L(x,v,t)=c(v)$，$c$ \emph{严格凸}，则最小化曲线具有\emph{常速度}。

  \textbf{证明（Jensen 不等式）：}
  \[
    A(\gamma)=\int_{0}^{1}c(\dot{\gamma}_{t})\dd t
    \;\ge\; c\!\left(\int_{0}^{1}\dot{\gamma}_{t}\dd t\right)
    = c(\gamma_{1}-\gamma_{0}).
  \]
  等号成立当且仅当 $\dot{\gamma}_{t}=\mathrm{const}$（常速直线）。$\square$

  \vspace{3pt}
  \deflab{例 7.4（Villani Example 7.4）。} 设 $M$ 为 Riemann 流形，取 $L(x,v,t)=|v|^{p}$，$p\ge1$：
  \begin{itemize}
    \item \textbf{$p>1$（严格凸）：} 最小化曲线是\TLterm{常速测地线}，满足
      $d(\gamma_{s},\gamma_{t})=|t-s|\cdot d(\gamma_{0},\gamma_{1})$，关联成本 $c=d^{p}$。
    \item \textbf{$p=1$（线性，非严格凸）：} 最小化曲线是测地曲线，允许任意重参数化。
  \end{itemize}

  \TLtakeaway{严格凸 Lagrange 量（$p>1$）通过 Jensen 不等式迫使最小化曲线等速；$c(x,y)=d(x,y)^p$ 是 Wasserstein 距离 $\Wp{p}$ 的成本函数。}
\end{frame}

% ── 度量空间中的测地线 ──────────────────────────────────────────────────────
\begin{frame}{度量空间中的测地线（例 7.8–7.9）}
  \deflab{例 7.8（Villani Example 7.8）。}
  设 $(\X,d)$ 为度量空间，曲线 $\gamma:[0,1]\to\X$ 的\TLterm{长度}为
  \[
    \mathrm{length}(\gamma):=\sup_{N}\sup_{0=t_{0}<\cdots<t_{N}=1}
    \sum_{i=0}^{N-1}d(\gamma_{t_{i}},\gamma_{t_{i+1}}).
  \]
  \TLterm{测地线}：长度等于端点距离的曲线（可变速，可重参数化为等速）。

  \deflab{例 7.9（Villani Example 7.9）。} 取 $L(x,v,t)=|v|^{p}$（$p>1$）。则

  $\gamma$ 为\TLterm{常速最小测地线} $\iff$ $d(\gamma_{s},\gamma_{t})=|t-s|\cdot d(\gamma_{0},\gamma_{1})$，$\forall s,t$.

  关联成本：$c(x,y)=d(x,y)^{p}$，最小化曲线是常速最小测地线。

  \deflab{测地长度空间（Villani p.298）。}
  若 $(\X,d)$ 中任意两点之间存在最短测地线，则称其为\TLterm{测地长度空间}。
  例：$\R^{n}$（直线段）、完备 Riemann 流形（Hopf–Rinow 定理）。

  \TLtakeaway{$p>1$ 的 $L^p$ 作用量选出常速最小测地线；测地长度空间保证最短路存在——二者共同构成位移插值理论的空间假设。}
\end{frame}

% ── 位移插值的构造 ────────────────────────────────────────────────────────────
\begin{frame}{位移插值的构造（一）：路径空间测度}
  \textbf{设置：} $(\X,d)$ 测地长度空间，$\mu_{0},\mu_{1}\in\Ppp{p}(\X)$，
  $\pi$ 为 $(\mu_{0},\mu_{1})$ 最优耦合（成本 $c=d^{p}$）。

  \textbf{第一步：} 对每对 $(x,y)$，选常速最小测地线 $\gamma^{x,y}:[0,1]\to\X$，
  满足 $d(\gamma^{x,y}_{s},\gamma^{x,y}_{t})=|t-s|\cdot d(x,y)$。

  \textbf{第二步：} 令 $e_{t}(\gamma):=\gamma(t)$，构造路径空间概率测度
  \[
    \Pi:=\int_{\X\times\X}\delta_{\gamma^{x,y}}\dd\pi(x,y).
  \]

  \textbf{第三步：} \TLterm{位移插值}为 $\mu_{t}:=(e_{t})_{\#}\Pi$。

  \TLtakeaway{位移插值 = 沿最优耦合对每对质量点取测地线后推前到 $t$ 时刻；$\Pi$ 是"随机测地线"的分布。}
\end{frame}

\begin{frame}{位移插值的构造（二）：$\R^{n}$ 中的显式公式}
  \textbf{特殊情形：} $\X=\R^{n}$，$c(x,y)=|x-y|^{2}$。当 $\mu_0\in\Pac(\R^n)$（绝对连续，见 §0）时，Brenier 定理（§2）保证最优映射 $T=\nabla\varphi$ 存在。

  直线段 $\gamma^{x,y}_{t}=(1-t)x+tT(x)$ 是常速最小测地线（$\R^n$ 中直线唯一），故
  \[
    \mu_{t}=\bigl((1-t)\Id+tT\bigr)_{\#}\mu_{0}.
  \]
  每质量点 $x$ 从 $x$ 以常速移向 $T(x)$，$t$ 时刻位于 $(1-t)x+tT(x)$。

  \begin{center}
  \begin{tikzpicture}[scale=1.0,>={Stealth[length=4pt]}]
    \fill[TLprimary] (0,0) circle (3pt) node[below,font=\scriptsize] {$x$};
    \fill[TLaccent]  (4,0.6) circle (3pt) node[below,font=\scriptsize] {$T(x)$};
    \draw[->,TLprimaryDark,thick] (0,0) -- (4,0.6);
    \fill[TLprimary!60!TLaccent] (2,0.3) circle (2.5pt) node[above,font=\scriptsize] {$\gamma_{1/2}^x$};
  \end{tikzpicture}
  \end{center}
  {\scriptsize 据 Villani 重绘。}

  \TLtakeaway{$\R^n$ 中的位移插值：每质量点沿直线匀速移动，$\mu_t=((1-t)\Id+tT)_\#\mu_0$；直线段唯一性保证构造的唯一性。}
\end{frame}

% ── 定义 7.19–7.20：动力学耦合 ──────────────────────────────────────────────
\begin{frame}{定义 7.19–7.20：动力学耦合与动力学最优耦合}
  \deflab{定义 7.19（动力学转移计划，Villani Def 7.19）。}
  \TLterm{动力学转移计划}是路径空间 $C([0,1];\X)$ 上的概率测度 $\Pi$。$\Pi$ 称为 $(\mu_{0},\mu_{1})$ 的\TLterm{动力学耦合}若
  $(e_{0})_{\#}\Pi=\mu_{0}$，$(e_{1})_{\#}\Pi=\mu_{1}$。

  \deflab{定义 7.20（动力学最优耦合，Villani Def 7.20）。}
  \TLterm{动力学最优转移计划}是集中在作用量最小化曲线全体 $\Gamma$ 上的概率测度 $\Pi$，使得
  $\pi_{0,1}:=(e_{0},e_{1})_{\#}\Pi$ 是 $(\mu_{0},\mu_{1})$ 的最优（静态）转移计划。

  等价地，$\Pi$ 是"随机作用量最小化曲线"的分布，端点构成最优耦合。

  \textbf{关系（路径空间 $\to$ 联合分布 $\to$ 边际）：}
  \begin{center}
  \begin{tikzpicture}[scale=0.85,>={Stealth[length=4pt]}]
    \node[draw=TLprimary,rounded corners=3pt,inner sep=5pt,font=\small] (A) at (0,0) {$\Pi\in\Pp(C([0,1];\X))$};
    \node[draw=TLaccent,rounded corners=3pt,inner sep=5pt,font=\small] (B) at (5,0) {$\pi_{0,1}=(e_0,e_1)_\#\Pi$};
    \draw[->,TLinkSoft,thick] (A) -- (B) node[midway,above,font=\scriptsize] {推前};
    \node[font=\scriptsize,TLinkSoft,below=6pt of A] {随机测地线};
    \node[font=\scriptsize,TLinkSoft,below=6pt of B] {最优静态耦合};
  \end{tikzpicture}
  \end{center}

  \TLtakeaway{$\Pi$（路径空间）$\to$ $\pi_{0,1}$（端点耦合，静态最优）$\to$ $\mu_t=(e_t)_\#\Pi$（位移插值）。}
\end{frame}

% ── 定理 7.21（陈述）────────────────────────────────────────────────────────
\begin{frame}{定理 7.21：位移插值（陈述）}
  \deflab{定理 7.21（位移插值，Villani Thm 7.21）。}

  设 $(\X,d)$ 为 Polish 空间，$(A)^{0,1}$ 为强制 Lagrange 作用量，成本 $c^{s,t}$ 连续，
  $C(\mu_{0},\mu_{1})<+\infty$。对连续路径 $(\mu_{t})_{0\le t\le1}$，以下三条\emph{等价}：
  \begin{enumerate}
    \item[(i)] $\mu_{t}=(e_{t})_{\#}\Pi$，其中 $\Pi$ 为某动力学最优耦合。
    \item[(ii)] 对任意 $t_{1}<t_{2}<t_{3}\in[0,1]$：
      $C^{t_{1},t_{2}}(\mu_{t_{1}},\mu_{t_{2}})+C^{t_{2},t_{3}}(\mu_{t_{2}},\mu_{t_{3}})
      =C^{t_{1},t_{3}}(\mu_{t_{1}},\mu_{t_{3}})$。
    \item[(iii)] $(\mu_{t})$ 是 $\Pp(\X)$ 上作用量泛函 $\mathcal{A}^{s,t}(\mu)$ 的最小化路径。
  \end{enumerate}
  满足上述条件的 $(\mu_{t})$ 称为\TLterm{位移插值}，至少存在一条。

  \TLtakeaway{定理 7.21 三条等价刻画：(i) 动力学最优耦合的边际，(ii) 测地型加法等式，(iii) 作用量最小化路径——三者互相等价。}
\end{frame}

% ── 推论 7.22（一）：陈述 ────────────────────────────────────────────────────
\begin{frame}{推论 7.22（一）：位移插值是 Wasserstein 测地线}
  \deflab{推论 7.22（Villani Corollary 7.22）。}

  设 $(\X,d)$ 为完备可分局部紧测地长度空间，$p>1$，$\mu_{0},\mu_{1}\in\Ppp{p}(\X)$。

  对连续路径 $(\mu_{t})_{0\le t\le1}\subset\Pp(\X)$，以下\emph{等价}：
  \begin{enumerate}
    \item[(i)] $\mu_{t}=(e_{t})_{\#}\Pi$，$\Pi$ 集中在常速最小测地线上且 $(e_{0},e_{1})_{\#}\Pi$ 为最优耦合。
    \item[(ii)] $(\mu_{t})_{0\le t\le1}$ 是 $\bigl(\Ppp{p}(\X),\Wp{p}\bigr)$ 中的\TLterm{常速测地线}：
      \[
        \Wp{p}(\mu_{s},\mu_{t})=|t-s|\cdot\Wp{p}(\mu_{0},\mu_{1}),\quad\forall s,t\in[0,1].
      \]
  \end{enumerate}
  对任意 $\mu_{0},\mu_{1}$，至少存在一条这样的测地线。

  \TLtakeaway{推论 7.22：Wasserstein 测地线 = "测地线的分布"（$\mu_t=\law(\gamma_t)$，$\gamma$ 为随机最优测地线）——联结最优传输与 Wasserstein 几何。}
\end{frame}

% ── 推论 7.22（二）：核心口号与意义 ─────────────────────────────────────────
\begin{frame}{推论 7.22（二）：核心口号与意义}
  \textbf{核心口号：}
  \[
    (\mu_{t})\text{ 是 }\Ppp{p}(\X)\text{ 中的测地线}
    \;\iff\;
    \mu_{t}=\law(\gamma_{t}),
  \]
  其中 $\gamma$ 是 $\X$ 中的随机常速最小测地线，端点构成最优耦合。

  \textbf{深层含义：}
  \begin{itemize}
    \item 最优传输（静态）$\longleftrightarrow$ Wasserstein 空间中的测地线（动态）。
    \item 推论 7.22 回答"$\Ppp{p}(\X)$ 中的直线是什么？"——就是位移插值！
    \item 后续章节（Ricci 曲率下界、位移凸性）以此为出发点。
  \end{itemize}

  \deflab{推论 7.23（唯一性，Villani Cor 7.23）。}
  若 (a) 最优传输计划唯一；(b) $\pi$-a.s. 两端点间最小化曲线唯一，则位移插值\emph{唯一}。

  \TLtakeaway{最优传输 $\iff$ Wasserstein 测地线：是定理 7.21 最深刻的几何含义，为后续章节奠定基础。}
\end{frame}

% ── 定理 7.21 证明思路 ────────────────────────────────────────────────────────
\begin{frame}{定理 7.21 证明思路（一）：三时刻引理}
  \textbf{目标。}\proofskipnote 证明 (ii) $\Rightarrow$ (i)：由三时刻加法等式构造路径空间上的动力学最优耦合 $\Pi$。
  \deflab{补一个事实：三时刻粘合。}
  若两个相邻最优耦合共享中间边际 $\mu_{t_{2}}$，则 Gluing Lemma（§0/§3）给出三元随机变量
  $(\gamma_{t_{1}},\gamma_{t_{2}},\gamma_{t_{3}})$，并保留两段边际耦合。
  \deflab{关键引理。}
  若 $(\gamma_{t_{1}},\gamma_{t_{2}})$ 与 $(\gamma_{t_{2}},\gamma_{t_{3}})$ 各自最优，且 (ii) 成立，则
  $(\gamma_{t_{1}},\gamma_{t_{3}})$ 也是最优耦合：
  \[
    \mathbb{E}\,c^{t_{1},t_{3}}(\gamma_{t_{1}},\gamma_{t_{3}})
    \le C^{t_{1},t_{2}}+C^{t_{2},t_{3}}
    \overset{\mathrm{(ii)}}{=}C^{t_{1},t_{3}}.
  \]
  左端又不小于最优值 $C^{t_{1},t_{3}}$，故只能处处等号，端点耦合最优。
  \deflab{补一个事实：限制性质。}
  {\small 同一条最小化曲线穿过中点时，它在每个子区间仍最小化；这是作用量最小化曲线的基本性质——最小化曲线的任意子段仍是最小化曲线（restriction/Markov property）。}
  \textbf{反向。} (i) $\Rightarrow$ (ii) 由 $\Pi$-a.s. 的最小化曲线直接积分计算得到。
  \TLtakeaway{三时刻引理把局部最优耦合和加法等式升级为跨端点最优；这是二进构造能够保持全局最优性的核心。}
\end{frame}

\begin{frame}{定理 7.21 证明思路（二）：二进构造与极限}
  \deflab{补一个事实：二进时刻。}
  二进时刻是 $j/2^{k}$；它们在 $[0,1]$ 中稠密，所以先控制这些时刻，再用连续性传到所有时刻。\proofskipnote
  \textbf{二进逼近（由 (ii) 构造 (i)）：}
  \begin{enumerate}
    \item 取 $(\mu_{0},\mu_{1})$ 的最优耦合 $\pi^{(0)}$。
    \item 在 $0,\tfrac{1}{2},1$ 上粘合两段最优耦合；上一帧引理保证端点耦合仍最优。
    \item 归纳到所有 $j/2^{k}$：反复粘合三时刻最优耦合，得到一致的离散随机路径。
    \item 在相邻二进点之间选择作用量最小化曲线，得到分段最小化路径的法律 $\Pi^{(k)}$。
  \end{enumerate}
  \deflab{本帧暂缓的技术点。}
  从 $\Pi^{(k)}$ 过渡到 $\Pi\in\Pp(C([0,1];\X))$ 需要测地线的可测选择，以及路径空间上的 tightness/compactness；Villani 第 7 章处理这些细节，主线只需二进边际一致且极限仍支撑在最小化曲线上。
  \TLtakeaway{(ii) $\Rightarrow$ (i) 的实质是：二进时刻反复粘合三时刻最优耦合，再用路径空间紧性取得连续随机最小化曲线。}
\end{frame}

% ── 计算例（一）：两个 Dirac 测度 ──────────────────────────────────────────
\begin{frame}{计算例（一）：两个 Dirac 测度的位移插值}
  \textbf{设置：} $\mu_{0}=\delta_{x}$，$\mu_{1}=\delta_{y}$，$x,y\in\X$，$p>1$。

  \textbf{最优耦合：} $\coup{\delta_{x}}{\delta_{y}}=\{\delta_{(x,y)}\}$（唯一），故 $\pi=\delta_{(x,y)}$，
  路径空间测度 $\Pi=\delta_{\gamma^{x,y}}$（唯一测地线上的 Dirac 测度）。

  \textbf{位移插值：}
  \[
    \mu_{t}=(e_{t})_{\#}\delta_{\gamma^{x,y}}=\delta_{\gamma^{x,y}_{t}}.
  \]
  中间分布仍是 Dirac 测度，集中在测地线 $\gamma^{x,y}$ 的 $t$ 时刻位置。

  \textbf{验证推论 7.22：}
  \[
    \Wp{p}(\mu_{s},\mu_{t})
    =\Wp{p}(\delta_{\gamma_{s}},\delta_{\gamma_{t}})
    =d(\gamma_{s},\gamma_{t})
    =(t-s)\,d(x,y)
    =(t-s)\,\Wp{p}(\mu_{0},\mu_{1}).
  \]
  常速测地线条件完美成立。

  \TLtakeaway{Dirac 测度的位移插值：中间分布仍为 Dirac 测度集中在测地线上——是定理 7.21 最简单的完整验证。}
\end{frame}

% ── 计算例（二）：一维高斯分布 ───────────────────────────────────────────────
\begin{frame}{计算例（二）：$\R$ 上两个高斯分布的位移插值}
  \textbf{设置：} $\mu_{0}=\mathcal{N}(m_{0},\sigma_{0}^{2})$，$\mu_{1}=\mathcal{N}(m_{1},\sigma_{1}^{2})$，$c(x,y)=(x-y)^{2}$。

  \textbf{最优映射：} 一维高斯的单调最优映射为仿射映射
  $T(x)=m_{1}+\tfrac{\sigma_{1}}{\sigma_{0}}(x-m_{0})$。

  \textbf{位移插值的显式公式：} $T_{t}:=(1-t)\Id+tT$，则
  \[
    \mu_{t}=(T_{t})_{\#}\mu_{0}=\mathcal{N}(m_{t},\sigma_{t}^{2}),
  \]
  其中
  \[
    m_{t}=(1-t)m_{0}+tm_{1},\qquad
    \sigma_{t}=(1-t)\sigma_{0}+t\sigma_{1}.
  \]

  \textbf{要点：} 中间分布仍是高斯！均值和标准差均线性插值，\emph{不出现双峰}。

  \textbf{对比线性插值：} $(1-t)\mathcal{N}(m_{0},\sigma_{0}^{2})+t\mathcal{N}(m_{1},\sigma_{1}^{2})$
  在 $m_{0}\ne m_{1}$ 时产生双峰高斯混合。

  \TLtakeaway{高斯的位移插值保持高斯性，均值和标准差均线性变化——与产生双峰的线性插值形成鲜明对比。}
\end{frame}

% ── 计算例（二）续：计算验证 ────────────────────────────────────────────────
\begin{frame}{计算例（二）续：高斯插值的验证}
  \textbf{验证 $\mu_{t}=\mathcal{N}(m_{t},\sigma_{t}^{2})$。}

  映射 $T_{t}(x)=(1-t)x+t\bigl(m_{1}+\tfrac{\sigma_{1}}{\sigma_{0}}(x-m_{0})\bigr)$ 整理为
  \[
    T_{t}(x)=\underbrace{\bigl(1-t+t\tfrac{\sigma_{1}}{\sigma_{0}}\bigr)}_{\alpha_{t}}x
             +\underbrace{t\bigl(m_{1}-\tfrac{\sigma_{1}}{\sigma_{0}}m_{0}\bigr)}_{\beta_{t}}.
  \]
  这是仿射映射，故 $(T_{t})_{\#}\mu_{0}$ 为正态分布。计算参数：
  \begin{align*}
    m_{t} &= T_{t}(m_{0})=(1-t)m_{0}+tm_{1},\\
    \sigma_{t}^{2} &= \alpha_{t}^{2}\,\sigma_{0}^{2}
    =\bigl((1-t)\sigma_{0}+t\sigma_{1}\bigr)^{2}.\quad\square
  \end{align*}

  \textbf{$W_{2}$ 距离：} $\Wtwo(\mu_{0},\mu_{1})^{2}=(m_{1}-m_{0})^{2}+(\sigma_{1}-\sigma_{0})^{2}$。

  \TLtakeaway{仿射映射的推前保持高斯性；方差公式 $\sigma_t=(1-t)\sigma_0+t\sigma_1$ 直接来自仿射斜率 $\alpha_t$。}
\end{frame}

% ── 计算例（三）：平移 ──────────────────────────────────────────────────────
\begin{frame}{计算例（三）：平移的位移插值}
  \textbf{设置：} $\mu_{1}=(\tau_{a})_{\#}\mu_{0}$，$\tau_{a}(x)=x+a$，$a\in\R^{n}$，$c=|\cdot|^{2}$。

  \textbf{最优映射：} $T(x)=x+a$（Brenier 定理：仿射映射为最优）。

  \textbf{位移插值：} $T_{t}(x)=x+ta$，故
  \[
    \mu_{t}=(\tau_{ta})_{\#}\mu_{0}.
  \]
  全体质量同步匀速平移距离 $ta$——整个分布形状不变，仅位置变化。

  \textbf{验证常速测地线条件：}
  \[
    \Wp{2}(\mu_{s},\mu_{t})
    =\Wp{2}\bigl((\tau_{sa})_{\#}\mu_{0},(\tau_{ta})_{\#}\mu_{0}\bigr)
    =|(t-s)a|
    =|t-s|\cdot\Wp{2}(\mu_{0},\mu_{1}).\quad\square
  \]

  \textbf{对比（与 Dirac 例子）：} 若 $\mu_{0}=\delta_{0}$，则 $\mu_{t}=\delta_{ta}$（Dirac 测度沿直线运动）；
  一般分布同样以同样速度整体平移，$\Wp{2}(\mu_{0},\mu_{1})=|a|$。

  \TLtakeaway{平移：整个分布匀速平移 $ta$，每质量点走直线——最直观的常速测地线验证，$\Wp{2}(\mu_0,\mu_1)=|a|$。}
\end{frame}

% ── Benamou–Brenier 前瞻 ─────────────────────────────────────────────────────
\begin{frame}{Benamou–Brenier 公式（前瞻）}
  \textbf{动力学视角（Benamou–Brenier，1999）。} 在 $\R^{n}$ 中，寻找速度场 $v$ 与路径 $(\mu_{t})$。

  \deflab{补一个事实：连续性方程。}
  $\partial_{t}\mu_{t}+\nabla\cdot(v_{t}\mu_{t})=0$ 表示质量守恒：速度场 $v_{t}$ 只搬运质量，不创造也不消灭质量。

  在此约束下最小化动能
  $\mathcal{A}(\mu,v)=\int_{0}^{1}\!\int_{\R^{n}}|v_{t}|^{2}\dd\mu_{t}\dd t$。

  \deflab{Benamou–Brenier 公式（Villani §7 文献注记）。}
  \[
    \Wtwo(\mu_{0},\mu_{1})^{2}
    =\min_{\partial_{t}\mu+\nabla\cdot(v\mu)=0}\mathcal{A}(\mu,v).
  \]

  \textbf{联系与意义。} 最优流就是位移插值路径
  $\mu_{t}=((1-t)\Id+tT)_{\#}\mu_{0}$；公式给出 $\Wtwo$ 的 PDE 表示、数值算法入口与 Otto 几何解释。

  \TLtakeaway{Benamou–Brenier：$W_2$ = 最小动能流；联结 PDE、数值计算与 Riemann 几何。}
\end{frame}

% ── §4 小结 ──────────────────────────────────────────────────────────────────
\begin{frame}{§4 小结：位移插值的核心结论}
  \begin{enumerate}
    \item \textbf{作用量框架：} $A(\gamma)=\int_{0}^{1}L(\gamma_{t},\dot{\gamma}_{t},t)\dd t$；
    严格凸 $L$（$p>1$）的最小化曲线为常速测地线（Jensen），关联成本 $c=d^{p}$。
    \item \textbf{构造：} 对最优耦合 $\pi$ 作测地线提升得 $\Pi$，$\mu_{t}=(e_{t})_{\#}\Pi$；
    $\R^{n}$ 中：$\mu_{t}=((1-t)\Id+tT)_{\#}\mu_{0}$。
    \item \textbf{定理 7.21：} 位移插值 $\iff$ 三时刻等号关系 $\iff$ 作用量最小化路径。
    \item \textbf{推论 7.22：} 位移插值 = $\Wp{p}$ 常速测地线：$\Wp{p}(\mu_{s},\mu_{t})=|t-s|\Wp{p}(\mu_{0},\mu_{1})$。
    \item \textbf{核心口号：} "概率测度空间的测地线 = 测地线的分布。"
  \end{enumerate}

  \textbf{下节预告：} §5 收缩性与位移凸性——位移插值的正则性和曲率下界。

  \TLtakeaway{位移插值统一了最优传输问题与 Wasserstein 空间的测地线几何，是整个 Part I 的几何核心。}
\end{frame}

% ── 习题 4 · Exercises（一）─────────────────────────────────────────────────
\begin{frame}{习题 4 · Exercises（一）}
  \diffI\ \textbf{习题 4.1.} 设 $\mu_{0}=\delta_{0}$，$\mu_{1}=\delta_{1}$，$\X=\R$，$p=2$。
  计算线性插值 $\mu_{t}^{\mathrm{lin}}=(1-t)\delta_{0}+t\delta_{1}$ 和位移插值 $\mu_{t}=\delta_{t}$ 之间的
  Wasserstein 距离 $\Wtwo(\mu_{t}^{\mathrm{lin}},\mu_{t})$。

  \hint{位移插值为 $\mu_{t}=\delta_{t}$；对于线性插值，利用 $\Wtwo^{2}((1-t)\delta_{0}+t\delta_{1},\,\delta_{t})$ 用最优匹配计算。}
  % SOLUTION: The only coupling of ((1-t)delta_0 + t*delta_1, delta_t) sends mass (1-t) from 0 to t and mass t from 1 to t. W_2^2 = (1-t)*t^2 + t*(1-t)^2 = t(1-t)(t+(1-t)) = t(1-t). So W_2 = sqrt(t(1-t)), maximized at t=1/2 with value 1/2.

  \vspace{6pt}
  \diffII\ \textbf{习题 4.2.} 设 $\mu_{0}=\mathcal{N}(0,1)$，$\mu_{1}=\mathcal{N}(0,4)$，$c(x,y)=(x-y)^{2}$。
  \begin{itemize}
    \item[(a)] 求最优传输映射 $T$。
    \item[(b)] 求位移插值 $\mu_{t}$ 的均值和方差。
    \item[(c)] 计算 $\Wtwo(\mu_{0},\mu_{1})$。
  \end{itemize}
  \hint{(a) $T(x)=(\sigma_{1}/\sigma_{0})x=2x$；(b) $m_{t}=0$，$\sigma_{t}=(1-t)\cdot1+t\cdot2=1+t$；(c) 由单调映射 $\Wtwo^{2}=\E[(X-T(X))^{2}]=\E[(X-2X)^{2}]=\E[X^{2}]=\sigma_{0}^{2}=1$，即 $\Wtwo=|\sigma_{1}-\sigma_{0}|=1$。}
  % SOLUTION: T(x)=2x; mu_t=N(0,(1+t)^2); W_2=|sigma_1-sigma_0|=|2-1|=1.

  \TLtakeaway{习题 4.1 揭示线性与位移插值的几何差距；习题 4.2 通过具体数字掌握一维高斯位移插值的全过程。}
\end{frame}

% ── 习题 4 · Exercises（二）────────────────────────────────────────────────
\begin{frame}{习题 4 · Exercises（二）}
  \diffIII\ \textbf{习题 4.3.} 设 $\mu_{0}$，$\mu_{1}$ 是 $\R^{n}$ 上均值为零、协方差矩阵分别为
  $\Sigma_{0}$，$\Sigma_{1}$ 的高斯分布。
  \begin{itemize}
    \item[(a)] 利用 Brenier 定理证明最优传输映射为线性映射 $T(x)=Ax$，其中
      \[
        A=\Sigma_{0}^{-1/2}\!\left(\Sigma_{0}^{1/2}\Sigma_{1}\Sigma_{0}^{1/2}\right)^{1/2}\!\Sigma_{0}^{-1/2}.
      \]
    \item[(b)] 写出位移插值 $\mu_{t}$ 的协方差矩阵 $\Sigma_{t}$。
    \item[(c)] 写出 $\Wtwo(\mu_{0},\mu_{1})^{2}$ 的迹公式。
  \end{itemize}
  \hint{(a) Brenier 定理：最优映射是凸函数的梯度 $T=\nabla\varphi$；高斯情形 $\varphi$ 是二次型，故 $T$ 线性且对称正定。(b) $\Sigma_{t}=((1-t)I+tA)\Sigma_{0}((1-t)I+tA)^{\top}$。(c) $\Wtwo^{2}=\mathrm{tr}(\Sigma_{0})+\mathrm{tr}(\Sigma_{1})-2\,\mathrm{tr}\bigl((\Sigma_{0}^{1/2}\Sigma_{1}\Sigma_{0}^{1/2})^{1/2}\bigr)$。}
  % SOLUTION: T(x)=Ax with A as given; mu_t=N(0,Sigma_t) with Sigma_t=((1-t)I+tA)Sigma_0((1-t)I+tA)^T; W_2^2=tr(Sigma_0)+tr(Sigma_1)-2tr((Sigma_0^{1/2}Sigma_1 Sigma_0^{1/2})^{1/2}).

  \TLtakeaway{高斯位移插值的矩阵公式是习题 4.2 的高维推广，训练对 Brenier 定理和矩阵平方根的运用。}
\end{frame}
